\begin{progprob}
    You are performing a massive study of Twitter data and have recorded over 1
    billion tweets. In your analysis, you frequently need to count the number
    of tweets posted after time $a$ but before time $b$. A linear
    time algorithm is not fast enough for your purposes.

    In a file named \python{logcount.py}, write a function named
    \python{logcount(arr, a, b)} which takes in a \emph{sorted} array \python{arr} and
    returns the number of elements of the array which are in the \emph{closed} interval
    \python{[a, b]}. You may assume that the elements of the array are numbers, but you
    should not assume that the elements are distinct (there may be duplicates).
    If the array is empty, you should return zero. If the interval $[a, b]$ is empty (as in the case
    when $a > b$, you should also return 0.
    To receive credit, your algorithm must have a worst case time complexity of $\Theta(\log n)$,
    where $n$ is the size of the array.

    Example: if \python{arr = [2, 3, 5, 5, 6, 9, 12]}, \python{a = 3.7}, and \python{b = 6.1},
    your function should return \python{3} since there are three elements within the closed
    interval $[3.7, 6.1]$; namely: 5, 5, 6.

    \textbf{Note}: This is our first \emph{programming problem} of the quarter. Programming
    problems will appear separate from the ``main'' homework on Gradescope and
    are (mostly) autograded. Upon submitting your code, make sure that the autograder passes
    \emph{all} of the tests -- it is checking to make sure that your file is named correctly,
    that the function runs, etc. However, passing all of these tests \textbf{does not} guarantee
    that you will receive full credit! After the homework due date, we will upload a more
    sophisticated autograder that will thoroughly check your code. A human grader will also
    quickly check your code to make sure that it adheres to any constraints imposed in
    the problem statement. 
    The only module your code may import is the \python{math} module.

    \begin{soln}
        We can use the general version of binary search we derived in discussion to find
        the first element $\geq a$ and the first element $>b$. The difference between
        the two indices is the length of the subarray that falls in $[a, b]$:

        \inputminted{python}{\thisdir/include-solution/logcount.py}
    \end{soln}
\end{progprob}
